\subsection{Functoriality of the Daniell extension}

The inclusion functors $\iota,\iota_{D},\iota_M$ are as in
Definition~\ref{def:inclusion-functors}.

\begin{lemma}\label{lem:daniell-ext-is-functor-1}
  Let $(X,L,I)$ and $(Y,M,J)$ be integration systems on sets $X$ and
  $Y$ respectively, and let $F:(X,L,I)\to(Y,M,J)$ be a morphism of integration
  systems.
  Then the map of base spaces $\underline{F}:Y\to X$ determines a morphism
  \[
    F^1:(X,L^1,I^1)\to(Y,M^1,J^1)
  \]
  in the category $\mathbf{DInt}$ of Daniell closed integration systems.
\end{lemma}
\begin{proof}
  Let $f\in L^1$ and let
  \[
    f\overset{\varphi_n}{\simeq}\sum f_n
  \]
  be a series representation of it.
  Since $F$ is a morphism of integration systems, for each $n$
  \[
    F^*f_n,\ F^*\varphi_n\in M .
  \]
  Moreover, for every $y\in Y$,
  \[
    |F^*f_n(y)|
    =|f_n(\underline{F}(y))|
    \le\varphi_n(\underline{F}(y))
    =F^*\varphi_n(y)
  \]
  and
  \[
    \sum_{n=1}^\infty J(F^*\varphi_n)
    =\sum_{n=1}^\infty I(\varphi_n)
    <+\infty.
  \]

  Furthermore, if $y\in Y$ satisfies
  \[
    \sum_{n=1}^\infty F^*\varphi_n(y)<+\infty
  \]
  then, putting $x:=\underline{F}(y)$,
  \[
    \sum_{n=1}^\infty\varphi_n(x)<+\infty .
  \]
  By the condition on the original series representation,
  \[
    \begin{aligned}
      F^*f(y)
      &=f(\underline{F}(y))\\
      &=\sum_{n=1}^\infty f_n(\underline{F}(y))\\
      &=\sum_{n=1}^\infty F^*f_n(y) .
    \end{aligned}
  \]
  Hence
  \[
    F^*f\overset{F^*\varphi_n}{\simeq}\sum F^*f_n
  \]
  is a series representation of $F^*f$, and $F^*f\in M^1$ follows. Therefore the
  map of base spaces $\underline{F}:Y\to X$ determines a morphism of Riesz
  systems
  \[
    F^1:(X,L^1)\to(Y,M^1) .
  \]

  Finally, by the definition of the extended integral via series
  representations and the fact that $F$ preserves the integral,
  \[
    \begin{aligned}
      J^1(F^*f)
      &=\sum_{n=1}^\infty J(F^*f_n)\\
      &=\sum_{n=1}^\infty I(f_n)\\
      &=I^1(f) .
    \end{aligned}
  \]
  Hence $F^1$ is a morphism of integration systems. Since the Daniell extension
  is a Daniell closed integration system by
  Lemma~\ref{lem:daniell-ext-is-complete}, the morphism $F^1$ is a morphism in
  $\mathbf{DInt}$.
\end{proof}

\begin{lemma}\label{lem:daniell-ext-is-functor-2}
  Assigning to an integration system $(X,L,I)$ on a set $X$ its Daniell
  extension $(X,L^1,I^1)$ determines a functor
  \[
    \mathcal{D}:\mathbf{Int}\to\mathbf{DInt} .
  \]
\end{lemma}
\begin{proof}
  It suffices to check the following.
  \begin{itemize}
    \item For an integration system $(X,L,I)$ on a set $X$, the
      identity morphism $\operatorname{id}_{(X,L,I)}$ satisfies
      $\mathcal{D}(\operatorname{id}_{(X,L,I)})=\operatorname{id}_{(X,L^1,I^1)}$.
    \item For integration systems $(X,L,I),(Y,M,J),(Z,N,K)$ on sets
      $X,Y,Z$ respectively and morphisms of integration systems
      $F:(X,L,I)\to(Y,M,J)$ and $G:(Y,M,J)\to(Z,N,K)$, we have
      $\mathcal{D}(G\circ F)=\mathcal{D}(G)\circ\mathcal{D}(F)$.
  \end{itemize}

  By Lemma~\ref{lem:daniell-ext-is-functor-1}, a morphism of integration systems
  \[
    F:(X,L,I)\to(Y,M,J)
  \]
  determines the morphism
  \[
    F^1:(X,L^1,I^1)\to(Y,M^1,J^1)
  \]
  induced by the same map of base spaces $\underline{F}:Y\to X$. We therefore
  set
  \[
    \mathcal{D}(F):=F^1 .
  \]

  First, the map of base spaces underlying the identity morphism
  $\operatorname{id}_{(X,L,I)}$ is $\operatorname{id}_X$. The morphism
  $\mathcal{D}(\operatorname{id}_{(X,L,I)})$ after the Daniell extension is
  induced by the same map $\operatorname{id}_X$, so
  \[
    \mathcal{D}(\operatorname{id}_{(X,L,I)})
    =\operatorname{id}_{(X,L^1,I^1)} .
  \]

  Next, take morphisms of integration systems
  \[
    F:(X,L,I)\to(Y,M,J),
    \qquad
    G:(Y,M,J)\to(Z,N,K) .
  \]
  The map of base spaces underlying the composite $G\circ F$ is
  \[
    \underline{G\circ F}=\underline{F}\circ\underline{G}:Z\to X .
  \]
  On the other hand, the maps of base spaces underlying
  $\mathcal{D}(F)=F^1$ and $\mathcal{D}(G)=G^1$ are $\underline{F}$ and
  $\underline{G}$ respectively, so the map of base spaces underlying
  $\mathcal{D}(G)\circ\mathcal{D}(F)$ is also
  $\underline{F}\circ\underline{G}$. Hence
  \[
    \mathcal{D}(G\circ F)
    =\mathcal{D}(G)\circ\mathcal{D}(F) .
  \]
  Therefore $\mathcal{D}$ is a functor.
\end{proof}
